\subsection{Uniform limit of the means and summation over levels} Regularity before stopping gives
\(p_i\le C_b/m\), a positive lower bound on diagonal vacancy, and a row--column gap~\cite[\S7]{pc}.
Consequently \(c_b I\le S_0+\Pi_m\le 4I\) for some \(c_b>0\). For the lower bound, the quadratic form is
the sum of \(\sum_i p_i(a_i^{\mathsf T}y)^2\) and \(\sum_f h_f y_f^2\), together with the gauge
projection term. The second term controls the diagonal components and the gap controls the axial
components orthogonal to the gauge. The upper bound follows because \(S_0\) has unit diagonal, its
off-diagonal mass per row is at most three, and the projection acts on its kernel.

Set \(A_m=S_0+\Pi_m-I\), \(M=(I+A_m)^{-1}=R+\Pi_m\), and \(r_b=1-c_b/4<1\). The polynomial
\(P_N=\tfrac14\sum_{h=0}^N[I-(I+A_m)/4]^h\) satisfies \(\|M-P_N\|_{\rm op}\le r_b^{N+1}/c_b\). To bound
an entry multiplied by \(m\), we use \(M=I-A_m+A_mMA_m\). Each row of \(A_m\) has absolute row sum at
most four and largest absolute entry \(O_b(m^{-1})\), so \(\|A_me_u\|_2=O_b(m^{-1/2})\) for every \(u\).
Hence \[
 \sup_{u,v}m\,|[A_m(M-P_N)A_m]_{uv}|\le C_b r_b^{N+1}.
\] This bounds the error in each entry of \(m(R-I)\) when \(M\) is replaced by \(P_N\) in the preceding
identity. Both the original and the truncated entries are bounded.

For fixed \(N\), powers of \(S_0-I\) are sums over sequences of adjacent lines. Insertions of \(\Pi_m\)
simplify by \((S_0-I)\Pi_m=-\Pi_m\). A product of walks is a finite sum over incidence types. If two
affine labels are identical in a type, the line is counted once. An additional coincidence imposes a
nonconstant equation and leaves \(O_{b,N}(n^{d-1})\) positions among the \(O(n^d)\) positions of a type
with \(d\) parameters. On the diagonal, a length-two return walk uses the same square in both steps. Its
contribution is \(3m\sum_{i\in u}p_i^2\), since there are three choices for the other line. Its excess
over the product of the means is precisely \(\theta_u=3\eta_u s_u-\int K_{uv}^2\,d\rho(v)\).
Intersections between diagonals are counted in the parity classes of Section~\ref{sec:contracciones}.

For each fixed type, we track the lines down to the auxiliary cut-off \(b/4\) as in~\cite[\S4]{tc}. At
fixed tolerance \(\eta>0\), outside an event of exponentially small probability,
\(mp_i=W(i)(1+O_b(\eta+m^{-1}))\) at live sites and \(h_f=u_f/U_f+O_b(\eta+m^{-1})\). For a fixed number
of marks, the one-step hazard is the sum of their individual hazards up to an error of
\(O_{b,N}(n^{-2})\) from double occupancies. Multiplying the survival factors down to level \(m\), an
axial mark contributes \(x\) and a diagonal mark contributes \(U_f(x)\), since \(\partial_x\log
U_f=(1-u_f)/U_f\). The approximation is uniform over the marks and over \(j\le m\le n\). Take
\(n\to\infty\) first and then \(\eta\downarrow0\). The normalized sums for each class are Riemann sums
over fixed polytopes. On each piece, the profile factors are continuous in \(x\) with a common modulus on
\([b,1]\). The truncated sums therefore converge uniformly over the levels to continuous walk integrals.

To apply the same argument to \(mV\), let \(\kappa_i=a_i^{\mathsf T}Ra_i-4\). The prefactor
expansion~\cite[\S7]{pc} gives the following formula, whose last sum runs over all lines. \[
 V=-14\sum_i p_i^3-\tfrac32\sum_i p_i^2\kappa_i
       +\tfrac94\sum_u\left(\sum_{i\in u}p_i^2\right)^2+O_b(m^{-2}).
\] The twelve off-diagonal entries in \(m\kappa_i\) contribute \(-12W(i)\) at the marked square. Together
with \(-14\) they give \(4\langle W^2\rangle_i\). The long walks and the corrected diagonal yield the
terms involving \(L_0,\theta\) in Section~\ref{sec:profile-data}. The square of the fiber sum gives the
last term there.

In the contractions the center has mass one, the exterior has mass at most one, and the four fibers
together have mass at most four. The free-line measure has mass \(\le6/b\). The square steps have total
mass \(\le8/b\). Expanding the difference of products multiplies each entry bound only by masses and
factors bounded at fixed \(b\). Thus the tail beyond degree \(N\) is \(\le C_b r_b^{N+1}\), uniformly in
the level for both the physical mean and its limiting integral. The insertion remainder at scale \(m^2\)
is \(O_b(m^{-1})\). Taking \(n\to\infty\) and then \(N\to\infty\) gives \begin{align*}
 \sup_{j<m\le n}\Big|
 \mathbb E_\nu\!\left[m^2\sum_{c\in i_m}p_c
       \big(P_c+(\delta_c^0)^2/2\big);H_j\right]-d(m/n)\Big|&\longrightarrow0,\\
 \sup_{j\le m\le n}\left|\mathbb E_\nu[mV(B_m);H_j]-\mathcal V(m/n)\right|
  &\longrightarrow0 .
\end{align*} Here \(i_m\) is the revealed row and \(H_j\) is the short regular arrival event. To average
over the row, first use the survival indicator before the row is chosen. Given the parent, the row is
uniform and the mean becomes \(m^{-1}\sum_c p_c(\cdot)\). Replacing that indicator by \(\mathbf1_{H_j}\)
for bounded observables costs \(O_b(\nu(H_j^c))\). The second formula uses all families in the potential.
This argument identifies a mean over parents and gives continuity of the limit on \([b,1]\). Numerical
evaluation of these integrals on grids also requires a bound on the spatial discretization error.

\paragraph{Endpoints and change of measure.} The uniform error \(\epsilon_n(b)\) in the first mean
contributes at most \(\epsilon_n(b)n\sum_{m>j}m^{-2}=o_b(1)\). Hence the \(P_c+(\delta_c^0)^2/2\) term
gives \(\int_b^1d(x)x^{-2}\,dx\). For \(\Delta_cV\), the conditional probability of a child given the row
under \(\nu\) is \(p_c(1+O_b(m^{-1}))\), and \(|\Delta_cV|\le C_bm^{-2}\). The substitution in the mean
costs \(O_b(m^{-3})\). The increments centered up to the stopping time form a martingale with second
moment \(O_b(\sum_{m>j}m^{-4})\). If \(M_\tau\) is their stopped sum, the differences are orthogonal in
\(L^2(\nu)\). Thus \(n|\mathbb E_\nu[M_\tau;H_j]|\le n(\mathbb E_\nu M_\tau^2)^{1/2}\nu(H_j)^{1/2}\).
Since \(\sum_{m>j}m^{-4}=O(j^{-3})\), the cost at scale \(n\) is \(O_b(n^{-1/2})\). On each history in
\(H_j\), the realized increments sum to \(V(B_j)-V(B_n)\). The second uniform mean gives \(\mathcal
V(b)/b-\mathcal V(1)\).

Up to the stopping time, \(n|A|\le C_b\) for the ordinary predictor. Lemma~\ref{lem:previa} shows that
\(n|R_Q(n)|\) is also bounded. On \(H_j\), set \(s=A-R_Q(n)+g_T+\Gamma\). There \(|s|=O_b(n^{-1})\), and
the inverse density on histories is \(e^{-s-g^\flat}\). By \eqref{eq:arrival-estimate}, \[
 \int_{H_j}|d\mu-d\nu|
 \le\mathbb E_\mu[|1-e^{-s}|+e^{-s}|1-e^{-g^\flat}|;H_j]=O_b(n^{-1}) .
\] The exit estimates in Appendix~\ref{app:signed} then give \(n\mathbb E_\mu A=n\mathbb
E_\nu[A;H_j]+o_b(1)\). For rooks, set \(u_e=g_T(M_e^{\rm fis})-g_T(M)\). The mask \(M_e^{\rm fis}\) is
obtained from the single minor \(M_e\) by applying the two diagonal exclusions of the child. The change
is \(O_b(m^{-2})\), including the step that causes exit. Thus \(\log\rho^\sharp-\log\rho=\log\mathbb
E_{\nu,i}e^{u_e}
 =\mathbb E_{\nu,i}u_e+O_b(m^{-4})\). The stopped density and the weighted exit bounds give
\(D(\mu^\tau\Vert\nu^\tau)=\mathbb E_\mu[A+\log X]-R_Q(n)=O_b(1)\). Let \(d_m\) be the conditional
relative entropies of the two child laws given the row, set to zero after stopping. The rows have the
same uniform law. The decomposition of relative entropy into its conditional increments gives
\(\sum_m\mathbb E_\mu
d_m=D(\mu^\tau\Vert\nu^\tau)\). Pinsker's inequality bounds the difference between the two means of
\(u_e\) by \(2C_bm^{-2}\sqrt{d_m/2}\). By Jensen's inequality and Cauchy--Schwarz, the sum of these
differences at scale \(n\) costs
\(O_b(n(\sum_{m>j}m^{-4})^{1/2}D(\mu^\tau\Vert\nu^\tau)^{1/2})=O_b(n^{-1/2})\). The realized rook
increments telescope, adding \((F(b)-1/12)/b-(F(1)-1/12)\) to the preceding endpoint terms. Thus at each
fixed cut-off, \[
 \lim_{n\to\infty} n\mathbb E_\mu A^\sharp
 =\int_b^1\frac{d(x)}{x^2}\,dx+\frac{G(b)}b-G(1)
 =\int_b^1\frac{\numerr(x)}{x^2}\,dx .
\]
